\section{Stealth Religions and Atheism}

As \textbf{Rene Guenon} often pointed out, there can be no conflict between science and religion as long as they remain within their respective domains. In particular, esoterism agrees with the findings of science. For example, science may claim that man has no free will, that he is subject to various unknown inner forces, and so on. Of course, esoterism asserts the same. The difference is that esoterism offers a way out: man must learn to become free, he must discover those hidden forces that direct his life. We recently discussed the esoteric nature of Darwinism. Now we will see if evolutionary biology offers us any further clues to human nature.

As the ``new atheist'' movement develops into a hardcore ideology, it begins to fragment into dissident groups. For our purposes, we can distinguish between two such groups:

\begin{itemize}


\item The \textbf{Brights}. They assert that science produces objectively true knowledge, and moreover, is the only source of knowledge. They interpret religious claims as equivalent to scientific claims about the empirical world. They then take an often hysterical ``moral'' stand against those claims under the principle that any belief that is not supported, or at least potentially supportable, by science as immoral.

\item \textbf{Evolutionary Biologists}. They take a more nuanced view.

Brains evolved to survive in the world, not to create objectively true theories. Furthermore, religions likewise evolved and persist because of their survival value. Hence, it would be incorrect to describe religions as simply immoral.
\end{itemize}


Since the brights are more widely known, philosophically unsophisticated, and very simple to understand -- it shouldn't take you more than 30 minutes or so to become a bright -- we will explore what the biologists say about it. That position is much more subtle, since it incorporates a wider range of biological phenomena. For example, the bright position is reductionist since it is overly focused on the ``gene''. However, genes alone can explain neither the organism nor its group behavior. \textbf{David Sloan Wilson} provides a useful summary of this position in This View of Life\footnote{\url{https://thisviewoflife.com/profile/david-sloan-wilson/}}.

\paragraph{The Argument from Reason}


First of all, Wilson undercuts the intellectual basis of the brights:

\begin{quotex}
As for the canons of rational thought, to the extent that brains evolved by natural selection, their main purpose is to cause organisms to behave adaptively in the real world--not to directly represent the real world.
\end{quotex}


Some of you may recognize this as a form of the Argument from Reason\footnote{\url{https://en.wikipedia.org/wiki/Argument_from_reason}} made famous by C S Lewis. The fundamental point is that if human reason or rationality arises from nonrational causes, then it has no validity of its own.

Wilson makes a less ambitious distinction between practical and factual realism. A belief is real in the practical sense if it leads to effective action, while a factually real belief corresponds to the external world. There is no question of an independent reason in this distinction. Not just religious, but also ideological and political, beliefs depend on what is practically rather than factually real. For example, what are called ``politically correct'' beliefs are often practically real, but not factual.

This is actually a very useful distinction and I'm certain that Gornahoor readers can find many examples. A dead giveaway is the repetition of simplistic slogans. I've read one prominent political leader who claimed that most people cannot understand something that won't fit on a placard.

Unfortunately, many debates end up in frustration since one side is committed to a practically real belief and the other to factually real beliefs. However, no quantity of facts will be persuasive in that exchange. Practical realism is usually accompanied with an intense emotional investment which would make a conversion to the factually real quite difficult psychologically.

Nevertheless, Wilson cannot resist making a normative claim:

\begin{quotex}
We need respect for factual realism as never before to arrive at practical solutions to life's complicated problems.
\end{quotex}


That is a good point with the proviso that esoterism considers the study of inner states of consciousness as fundamental to factual realism. This is what Julius Evola called ``metaphysical positivism''.

Another point is that Wilson's claim has a goal: viz., ``practical solutions''. Yet, this gives us pause since most people are committed to practical realism. How would the factual realist motivate the practical realist, if not with the ``noble lie''?

\paragraph{Stealth Religions}


Wilson defines a stealth religion as a factually incorrect belief system that functions like a religion. Specifically, its main purpose is not to describe the main world but rather to motivate a suite of behaviours. The difference is that a stealth religion does not accept the existence of anything transcendent.

Wilson asserts that the real world involves ``messy tradeoffs''; for example, the effects of an action may be beneficial, harmful, or a mix of both for the parties involved. A stealth religion evades those tradeoffs by turning them into absolutes, either good or bad for everyone, anytime.

Now even animals know facts, but what is truly human is the intellectual capacity to grasp the concepts, or eternal ideas, that explain the facts. That is the definition of rationality. The argument from reason denies, then, that rationality per se is a ``natural process''. That puts Wilson in an intellectual bind. If legitimate religions are the result of evolutionary forces, then so are the stealth religions. But which intellectual system is not a stealth religion by his definition?

Certainly multiple political theories cannot all be true, i.e., all, or in the best case, all but one, of them are stealth religions because they arose arbitrarily and must be factually incorrect in part or in whole. Just a quick survey of the news today reveals slogans like, ``I am a progressive who gets things done,'' without specifying any thing in particular. Or ``People will respond to the message of hope,'' but hope for what? Nevertheless, they persist, since they are reflect group cohesion, as we shall see later.

So either reason is transcendent to the world, or, as a natural process, it can only create competing and incompatible belief systems. This, however, is not the lesson that the atheists will learn from Wilson.

\paragraph{The Proud Atheist}


Wilson defends himself from being a stealth religionist by claiming to be a ``proud'' atheist. I have no idea what pride has to do with it. In particular, he is proud because he does not believe in people sitting on clouds intervening in natural processes. If that is his actual view on religious beliefs, then it is impossible to take him seriously as a scientist competent to explain religious phenomena.

Nevertheless, Wilson is even-handed as he points out how pseudo-sciences like the New Atheism, strident forms of environmentalism, and probably other beliefs that he doesn't dare express, act as stealth religions. By pseudo-science, I mean a belief system that claims to be based strictly on science but is actually a stealth religion.\footnote{Wilson's understanding of God is no better than that of folk religion. However, there are more sophisticated versions of folk religion that go by the name of ``theistic personalism''. I would agree with Wilson were that the only alternative. However, the Traditional understanding, sometimes called ``classical theism'', is the only correct one. For a brief discussion of this, please listen to this lecture by David Bentley Hart: \url{https://www.youtube.com/watch?v=ve1s4amWCg4}.}

\paragraph{Proximate and Ultimate Causation}


Wilson makes a distinction between ultimate and proximate causation is evolutionary theory. This is what he means by them, although this is not his own explanation. Neo-Darwinism recognizes two primary factors: one is the characteristics of the organism and the other is the environment which acts as the selector, choosing the winners among the organisms. Both causes are necessary to understand ``evolution by natural selection''.

\begin{itemize}


\item \textbf{Ultimate Causation}. This refers to the environmental factors. For example, flowers bloom in spring because it is the optimal time. (Any earlier, then frost might kill it, later, then the autumn frost may thwart its development.)

\item \textbf{Proximate Causation}. This refers to the organism. To complete the example, the flower has a physiological mechanism that causes it to bloom in the spring.
\end{itemize}


Wilson then notes the similarity of these concepts with the vertical and horizontal dimensions of religion. The vertical dimension refers to the believers' relationship to God or the transcendent, and the horizontal dimension refers to their relationship to each other, i.e., the community.

\paragraph{Religion as a Natural Phenomena}


Wilson proposes six possible naturalistic explanations, of uneven plausibility, for the origin of religions. He more or less likes them all, but in the end settles on the explanation he calls the ``\textbf{Superorganism Hypothesis}''.

\begin{quotex}
If you could say only one thing about religion, it would be this: Most enduring religions have what Emile Durkheim called ``secular utility.'' They define, motivate, and coordinate groups to achieve collective goals in this life. They promote cooperation within the group and bristle with defenses against the all-important problem of cheating. Using the terms that I introduced in part I, they score high on practical realism, no matter how much they depart from factual realism along the way.
\end{quotex}


So how exactly does the individual benefit? Wilson dares to tell us:

\begin{quotex}
With respect to the individual benefits of religion, suppose that you discover a grand mansion, better than anything that you could have constructed on your own, with a sign on the door that says ``Welcome! Move right in!'' You would be a fool to refuse, and your decision might be purely selfish, with only your own welfare in mind.
\end{quotex}


Nevertheless, Wilson decides to be a fool, but a ``proud'' fool.

So the group forms a superorganism that motivates group behavior in beneficial ways, even if its empirical claims are not factually real. The hypothesis explains the horizontal dimension of religion, but Wilson conveniently forgets the vertical dimension. The real problem is specificity: why this religion and not another?

It is one thing to claim that a particular religion arose due to natural processes, but how did it take a particular form? That can only be explained by the vertical dimension, i.e., a revelation from above.\footnote{We will say more about the ``superorganism'' in the near future, based on Vladimir Solovyov's lecture on August Comte.}

\paragraph{Religion as Factually Real}


Scientists tend to approach religion as if it were a scientific theory that tries to explain empirical facts. In that respect it must seem to be absurd. That may be true for folk religions, but not for the deeper understanding of religious texts. Gregory of Nyssa explains:

\begin{quotex}
If one does not read scripture in a ``philosophical'' fashion, one will see only myths and contradictions.
\end{quotex}


This is the Protestant legacy with its sola scriptura doctrine understood literally. The cultured despisers of religion then rely on the literal interpretation rather than the philosophical, or allegorical, interpretation. For more on this, you can listen to this lecture by David Bentley Hart\footnote{\url{https://www.youtube.com/watch?v=5QPNvri-zxs}}.

\paragraph{Noble Lie}


Wilson is in an odd position. If, as a scientist, he realizes that a religion promotes group identity and mutual cooperation, not to mention individual benefits, then is he morally obliged to remain silent about its real (according to him) origins? What if the promotion of atheism brings about social turmoil and decadence? How can he know that a group commitment to atheism will serve as a new social glue? Given the time spans involved in evolutionary theories, it would take generations for the full effects of a belief system to unfold.

In other words, if all opinions are arbitrary and bigoted, then only a Noble Lie will enforce one of them over the others.


\begin{center}* * *\end{center}

\begin{footnotesize}
\begin{sffamily}
\texttt{Wayne Ferguson on 2016-02-12 at 12:02 said:}

The resurrection can be literally true without the Easter narratives being factual in a historical sense (tombs opening during the crucifixion, earthquakes, angels appearing in the garden, Jesus levitating into the clouds, etc.). While there may not be any down-side to believing that these things happened as reported-- as long as such beliefs are honestly held and not simply affirmed out of a fear or a desire to conform --we can remove a stumbling block to living faith (in the minds of many) if we can acknowledge that these narratives (as with many of the miracles) can also be thought of as symbolic (i.e. symbols woven into a quasi-historical narrative that point us to the literal/spiritual truth of the resurrection and the life). ``The Spirit of the Lord is upon me, because he has anointed me to bring good news to the poor. He has sent me to proclaim release to the captives and recovery of sight to the blind, to let the oppressed go free, to proclaim the year of the Lord's favor'' (Luke 4:18-19). ``Therefore he says, Awake you that sleep, and arise from the dead, and Christ shall give you light'' (Ephesians 5:14). Blind eyes continue to be opened today... Demons are exercised... Spiritually and psychologically lame rise up and walk... The dead are raised... Miracles have not ceased to happen, we have just become fixated on a particular idea of what a miracle is. That idea had a lot of rhetorical power at one time, but it seems to me that it has (on balance) become a stumbling block today.

\hfill

\texttt{White Rabbit on 2016-02-11 at 20:55 said:}

``Is the intelligent atheist in fact obliged to be very quiet about his atheism?'' Only if that thought randomly happens to possess him... otherwise his passion for discourse and ``truth'' will ``oblige'' him to write some more!

\hfill

\texttt{Mark Citadel on 2016-02-11 at 18:54 said:}

KJ, from my own understanding you must learn to both read a text as broadly differentiated as the Bible in a correct way, and understand the context of each section. The Resurrection I think is necessarily literally true (given as a testimony) whereas something like the literal Adam and Eve one doesn't have to be as stringently committed to, since their description lies beyond he fog, back into the age of myth, where symbolism was hardly discernible from the real. ``Sloan defines a stealth religion as a factually incorrect belief system that functions like a religion. Specifically, its main purpose is not to describe the main world but rather to motivate a suite of behaviours. The difference is that a stealth religion does not accept the existence of anything transcendent.'' He just defined the Cult of Progress. The last question you raise is certainly interesting. Is the intelligent atheist in fact obliged to be very quiet about his atheism? hmmm... that will require some pondering.

\hfill

\texttt{Cologero on 2016-02-11 at 07:38 said:}

@Kj, I'm just repeating the traditional understanding. The Bible cannot be read simply as an historical or scientific text. It has been read on four levels, of which the literal or physical interpretation is the first level.

\hfill

\texttt{KJ on 2016-02-11 at 04:49 said:}

Cologero, Are you claiming that *all* of the Bible should be read allegorically? Do you believe any of it to be literally true? Whatever about his involvement in esoteric Christianity there can hardly be any doubt that Dante and his fellow Christians of the High Middle Ages believed the resurrection to be literally correct.

\end{sffamily}
\end{footnotesize}


